% TOP_PROBLEM: {"id": 2629, "title": "Kourovka Notebook Problem 21.120", "category_id": 17, "category": {"id": 17, "name": "group_theory", "display_name": "Group Theory", "description": "Problems about groups, group actions, representations, and related algebraic structures.", "slug": "group-theory", "order_index": 17, "created_at": "2026-07-31T15:26:25.671Z"}, "status": "open", "problem_number": "KOU-21.120", "set_id": 10}
% FIRST_POSTED: 2026-10-01
% ENTRY_KIND: statement_only
% UnsolvedMath ID 2629; source catalogue code KOU-21.120
% !TeX program = lualatex
% Definition and sources only; this file is not an LLM solution attempt.
\documentclass[11pt,a4paper]{article}
\usepackage[margin=25mm]{geometry}
\usepackage{fontspec}
\setmainfont{lmroman10-regular.otf}[BoldFont=lmroman10-bold.otf,
 ItalicFont=lmroman10-italic.otf,BoldItalicFont=lmroman10-bolditalic.otf]
\usepackage{amsmath,amssymb,mathtools}
\usepackage{unicode-math}
\setmathfont{latinmodern-math.otf}
\AtBeginDocument{\let\setminus\smallsetminus}
\usepackage{xurl,hyperref}
\hypersetup{colorlinks=true,urlcolor=blue}
\setcounter{secnumdepth}{0}
\setlength{\parindent}{0pt}
\setlength{\parskip}{5pt}
\setlength{\emergencystretch}{3em}
\begin{document}
\section*{Kourovka Notebook Problem 21.120}\label{2629}
{\small UnsolvedMath ID 2629 · KOU-21.120 · Group Theory · Kourovka Notebook - New Problems, Issue 21}
\subsection{Definitions and mathematical statement}
Fix a prime $p$. A pro-$p$ group is an inverse limit of finite $p$-groups, with its inverse-limit topology. The pro-$p$ completion of an abstract group $H$ is $\widehat H_p=\varprojlim H/N$, where $N$ runs over normal subgroups of finite $p$-power index.

A pro-$p$ group is strictly finitely presented if it is isomorphic to $\widehat H_p$ for some finitely presented abstract group $H$. In the relative version, $H$ need only be finitely presented within a finitely based variety of groups. Such a variety is specified by finitely many group identities. A finite presentation within it means a quotient of a relatively free group on finitely many generators by the normal closure of finitely many relators.

A pro-$p$ group $G$ is finitely axiomatizable among pro-$p$ groups if there is one parameter-free first-order sentence $\sigma$, using multiplication, inverse, identity, and equality, such that $G\models\sigma$ and every pro-$p$ group satisfying $\sigma$ is isomorphic to $G$ as a topological group. Quantifiers in the sentence range over elements, not over open subsets.

\textbf{(a)} Does there exist a strictly finitely presented pro-$p$ group that is not finitely axiomatizable in this class? Ask the same question with ``strictly'' replaced by ``relatively strictly''.

\textbf{(b)} Is every finitely generated free pro-$p$ group finitely axiomatizable? The free pro-$p$ group on $d$ generators is the pro-$p$ completion of the abstract free group of rank $d$; maps from its generators into any pro-$p$ group extend uniquely to continuous homomorphisms.
\subsection{Short English statement}
Can strict or relative finite presentability fail to give a single first-order characterization among pro-$p$ groups, and are free pro-$p$ groups finitely axiomatizable?
\subsection{Sources}
\begin{itemize}
\item[S1] ULAM.AI and the UnsolvedMath Contributors, supplied problems.json, record 2629 (KOU-21.120), Kourovka Notebook - New Problems, Issue 21. Catalogue identity and source transcription; CC BY 4.0.
\url{https://huggingface.co/datasets/ulamai/UnsolvedMath}.
\item[S2] The Kourovka Notebook, 21st edition, new problems, Problem 21.120. Expanded formulation of the supplied catalogue statement.
\url{https://arxiv.org/pdf/1401.0300v47}.
\item[S3] D. Segal, On some relatively free pro-p groups; framework for relative presentations and finite axiomatizability.
\url{https://arxiv.org/abs/2505.04816}.
\end{itemize}
\end{document}
